\documentclass[11pt,a4paper]{article}

\usepackage[utf8]{inputenc}
\usepackage[T1]{fontenc}
\usepackage{lmodern}
\usepackage[french]{babel}
\usepackage[margin=2.3cm]{geometry}
\usepackage{microtype}
\usepackage{booktabs}
\usepackage{array}
\usepackage[table]{xcolor}
\usepackage{amsmath}
\usepackage{amssymb}
\usepackage{mathtools}
\usepackage{enumitem}
\usepackage[most]{tcolorbox}
\usepackage{tabularx}
\usepackage{eurosym}
\usepackage{caption}
\usepackage{titlesec}
\usepackage{hyperref}

\definecolor{navy}{HTML}{1F3864}
\definecolor{bluemid}{HTML}{2E5496}
\definecolor{lightblue}{HTML}{D6E0F0}
\definecolor{accent}{HTML}{C0491F}
\definecolor{green}{HTML}{2E6B4F}
\definecolor{lightgreen}{HTML}{D6EADF}

\hypersetup{colorlinks=true, linkcolor=bluemid, urlcolor=bluemid,
  pdftitle={Conformite multi-etats et trajectoire du SCR DORA},
  pdfauthor={Kelian Kaddouri}}

\newtcolorbox{cle}[1][]{colback=lightblue!40,colframe=navy,boxrule=0.6pt,
  arc=2pt,left=8pt,right=8pt,top=6pt,bottom=6pt,before skip=9pt,after skip=9pt,#1}
\newtcolorbox{portee}[1][]{colback=lightgreen!50,colframe=green,boxrule=0.6pt,
  arc=2pt,left=8pt,right=8pt,top=6pt,bottom=6pt,before skip=9pt,after skip=9pt,#1}

\titleformat{\section}{\normalfont\large\bfseries\color{navy}}{\thesection}{0.6em}{}
\titlespacing*{\section}{0pt}{2.0ex plus .3ex minus .2ex}{1.0ex plus .2ex}
\titleformat{\subsection}{\normalfont\normalsize\bfseries\color{bluemid}}{\thesubsection}{0.6em}{}
\captionsetup{font=small,labelfont={bf,color=navy},labelsep=period,skip=6pt}

\title{\vspace{-1.2cm}\color{navy}\bfseries Conformité multi-états et trajectoire du SCR DORA\\[2pt]
  \large Du Vasicek dirigé à un capital qui dépend de l'état de conformité, par pilier et dans le temps}
\author{Kélian Kaddouri}
\date{\today}

\begin{document}
\maketitle
\vspace{-0.5cm}

\begin{cle}
\textbf{Idée.} Le chantier cascade (Vasicek dirigé) chiffre un SCR pour une entité
\emph{figée}. On ajoute ici que ce SCR \textbf{dépend de l'état de conformité} de l'entité,
état qui n'est ni binaire ni statique. Trois briques : (i)~un Vasicek \textbf{multi-états} qui
définit les états de conformité par des seuils ordonnés sur une latente et les fait piloter
la cascade~; (ii)~une déclinaison \textbf{par pilier} qui décompose le surcoût de
non-conformité~; (iii)~une chaîne de \textbf{Markov} qui donne la \textbf{trajectoire} du
capital le long de la mise en conformité. C'est le remplacement dynamique de la variable
latente statique de l'ancien mémoire.
\end{cle}

\noindent Ce document accompagne les scripts \texttt{16} à \texttt{20b} du dossier
\texttt{vasicek\_lab} (table en fin de document). Il est complété au fil du chantier.

\paragraph{Deux décisions de cadrage (validées).}
\begin{enumerate}[leftmargin=1.5em,itemsep=1pt,topsep=2pt]
\item \textbf{Remplacement.} La chaîne multi-états remplace la variable latente de conformité
  statique du mémoire. Un seul modèle de conformité, pas deux concurrents.
\item \textbf{Par pilier.} L'état de conformité est propre à chaque pilier, ce qui fait
  interagir la conformité avec la cascade (un pilier non conforme propage plus et détecte
  moins que ses voisins). On l'établit d'abord en état global (\S\ref{sec:global}), cas
  particulier du par-pilier (\S\ref{sec:pilier}).
\end{enumerate}

\section{Le Vasicek multi-états : état global}
\label{sec:global}

\subsection{Les états viennent de seuils ordonnés sur une latente (1.1)}

Chaque entité (\S\ref{sec:global}) ou chaque pilier (\S\ref{sec:pilier}) porte une
\textbf{capacité de conformité latente} $C^\star\sim\mathcal N(0,1)$, décomposée à la Vasicek
en un facteur systémique commun et un choc idiosyncratique :
\begin{equation}
C^\star \;=\; \gamma\,\Theta \;+\; \sqrt{1-\gamma^2}\,\varepsilon,
\qquad \Theta,\varepsilon\sim\mathcal N(0,1),\quad \gamma=0{,}68.
\end{equation}
Au lieu d'un seul seuil de défaut, on en pose \textbf{deux}, ordonnés, qui découpent trois
états de conformité :
\begin{equation}
\text{NC si } C^\star\le K_{\text{bas}}, \qquad
\text{PC si } K_{\text{bas}}<C^\star\le K_{\text{haut}}, \qquad
\text{C si } C^\star>K_{\text{haut}}.
\end{equation}
Les seuils sont fixés par les probabilités d'état \textbf{ancrées} (enquêtes de conformité
fin 2025) : $K_{\text{bas}}=\Phi^{-1}(p_{\text{NC}})$ et
$K_{\text{haut}}=\Phi^{-1}(p_{\text{NC}}+p_{\text{PC}})$, avec
$(p_{\text{NC}},p_{\text{PC}},p_{\text{C}})=(0{,}35\,;0{,}35\,;0{,}30)$, d'où
$K_{\text{bas}}=-0{,}385$ et $K_{\text{haut}}=+0{,}524$.

\paragraph{Le facteur systémique rend les probabilités conditionnelles.} La probabilité
d'être dans chaque état sachant l'environnement $\Theta$ s'écrit
\begin{equation}
\mathbb P(\text{NC}\mid\Theta)=\Phi\!\Big(\tfrac{K_{\text{bas}}-\gamma\Theta}{\sqrt{1-\gamma^2}}\Big),
\qquad
\mathbb P(\text{C}\mid\Theta)=1-\Phi\!\Big(\tfrac{K_{\text{haut}}-\gamma\Theta}{\sqrt{1-\gamma^2}}\Big).
\end{equation}
Intégrée sur $\Theta$, la marginale redonne exactement l'ancrage (contrôle de cohérence). En
crise systémique ($\Theta=-2{,}5$), la masse bascule massivement vers la non-conformité :
$\mathbb P(\text{NC})$ passe de $35\,\%$ à $96\,\%$. C'est la version à trois états ordonnés de
l'amplification systémique que portait l'ancienne variable latente.

\subsection{L'état pilote les trois canaux (1.2)}

L'état de conformité $c$ agit sur la cascade par trois canaux, câblés sans réécrire
l'agrégation (le moteur euro-cascade est réutilisé tel quel) :
\begin{enumerate}[leftmargin=1.5em,itemsep=1pt,topsep=2pt]
\item \textbf{Fréquence} $\lambda$ : multiplicateurs de scénario $S_0/S_1/S_2$ sourcés
  (ENISA, Ponemon, Microsoft Research). Un état dégradé produit plus d'incidents.
\item \textbf{Propagation} $g$ : une entité plus conforme contient mieux la contagion, donc
  un gain plus faible. La propension à propager depuis un pilier est $e_j=g(c_j)\,s_j/\max_k s_k$
  (gain du pilier \textbf{source}).
\item \textbf{Détection} $p_u$ : le taux de matérialité. Une entité conforme intercepte et
  remédie tôt, donc une moindre fraction d'incidents atteint la sévérité matérielle.
\end{enumerate}
On présente trois lectures emboîtées : \textbf{A}~fréquence seule (comparable au mémoire),
\textbf{B}~$+$~propagation, \textbf{C}~$+$~détection. Cet emboîtement rend visible la
contribution marginale de chaque canal et évite qu'un canal domine en silence.

\subsection{Résultats : SCR par état et Delta\_DORA}

Le SCR de chaque état est la VaR à $99{,}5\,\%$ de la perte annuelle agrégée par la cascade.
Le \textbf{Delta\_DORA} est le surcoût strictement imputable à la non-conformité, à graine
Monte-Carlo \textbf{commune} entre états (l'écart ne vient que du changement d'état) :
\begin{equation}
\Delta_{\text{DORA}}(\text{état}) \;=\; \mathrm{SCR}(\text{état}) - \mathrm{SCR}(\text{C}).
\end{equation}

\begin{table}[h]\centering\small
\renewcommand{\arraystretch}{1.25}
\begin{tabular}{@{}l ccc@{}}
\toprule
SCR (M\euro, OpRisk) & lecture A & lecture B & lecture C \\
\midrule
Conforme               & 8301  & 6664  & 5932 \\
Partiellement conforme & 11059 & 9625  & 9625 \\
Non conforme           & 15074 & 15074 & 16595 \\
\bottomrule
\end{tabular}
\caption{SCR par état de conformité, entité type du mémoire (script~\texttt{16}). Échelle
monotone C~$<$~PC~$<$~NC. L'état intermédiaire PC est l'apport par rapport au mémoire.}
\end{table}

Le Delta\_DORA (NC contre C), en bootstrap deux niveaux (multiplicateurs $+$ sévérité, sévérité
de queue OpRisk ré-échantillonnée sur les $91$ excès réels comme au mémoire, non tirée dans une
approximation normale), reste du même ordre que le mémoire, amplifié par la propagation : OpRisk
médiane $7401$~M\euro{} (IC90 $[2830\,;38793]$, mémoire $3879$), PRC médiane $3668$~M\euro{}
(IC90 $[3201\,;4094]$, mémoire $2015$), et $100\,\%$ des tirages positifs. Le surcoût
intermédiaire (PC contre C) est chiffré pour la première fois : OpRisk médiane $2556$~M\euro{}
(IC90 $[963\,;12768]$), PRC $1257$~M\euro{} (IC90 $[1077\,;1466]$). L'intervalle OpRisk reste
large parce que la queue l'est (facteur $\sim 14$) : la direction du surcoût est certaine, son
niveau non pinçable, exactement la thèse du mémoire.

Le pendant PRC de cette honnêteté est le script~\texttt{21} : la sévérité PRC étant dérivée
des enregistrements compromis (conversion log-log de Jacobs 2014), on réinjecte dans le
bootstrap le résidu de la régression d'origine ($\mathrm{RSE} = 0{,}523$ en log naturel,
$R^2 = 0{,}51$). L'IC ne se rouvre presque pas (facteur $\approx 1{,}4$ avec comme sans
résidu : $15\,053$ enregistrements moyennent le bruit dans l'ajustement GPD et le cap borne
la queue), mais le niveau monte de $7\,\%$ (médiane NC contre C de $3\,600$ à
$3\,861$~M\euro, queue dérivée $\xi$ de $1{,}03$ à $1{,}07$). L'étroitesse de l'IC PRC n'est
donc pas un artefact de la conversion déterministe~; l'incertitude restante est celle des
coefficients $(a,b)$, à traiter en axe de sensibilité.

\section{État de conformité par pilier}
\label{sec:pilier}

Chaque pilier porte désormais son propre état. Les trois canaux deviennent indexés par pilier
(dictionnaires $\lambda_j$, $g_j$, $p_{u,j}$), passés à un moteur \texttt{simulate\_euro\_pp}
qui se réduit exactement à la version globale quand tous les piliers partagent un état
(contrôle de cohérence : configurations homogènes redonnent le~\S\ref{sec:global}).

\begin{cle}
\textbf{Décomposition du Delta\_DORA.} On bascule \textbf{un seul} pilier de Conforme à Non
conforme, les autres restant conformes, et l'on mesure le surcoût propre du pilier $k$ :
\[
\Delta_k \;=\; \mathrm{SCR}(k~\text{seul NC, autres C}) - \mathrm{SCR}(\text{tous C}).
\]
\end{cle}

\begin{table}[h]\centering\small
\renewcommand{\arraystretch}{1.25}
\begin{tabular}{@{}l cc@{}}
\toprule
Pilier basculé en NC (OpRisk) & $\Delta_k$ (M\euro) & part \\
\midrule
P1 gouvernance   & 3481 & 37\,\% \\
P4 tiers ICT     & 2658 & 28\,\% \\
P2 incidents     & 1634 & 17\,\% \\
P3 tests         & 1094 & 11\,\% \\
P5 partage       & 665  & 7\,\% \\
\bottomrule
\end{tabular}
\caption{Contribution de chaque pilier à la non-conformité (script~\texttt{16b}). Le classement
$P1>P4>P2>P3>P5$ \textbf{reproduit le classement ROOT} du modèle qualitatif (validation
croisée euro~$\leftrightarrow$~ordinal).}
\end{table}

\noindent Deux enseignements. D'abord, la \textbf{validation croisée} : la hiérarchie des
piliers obtenue par le chiffrage en euros coïncide avec le classement ROOT posé par jugement
d'expert dans le modèle qualitatif, deux voies indépendantes qui convergent. Ensuite,
\textbf{l'interaction} : le Delta total (tous NC contre tous~C) ne coïncide pas avec la somme
des $\Delta_k$, précisément parce que la cascade fait interagir les piliers (un pilier non
conforme amplifie aussi les contagions amorcées ailleurs qui le traversent), ce qu'un modèle
sans propagation ne peut pas capturer. Sa mesure demande de la prudence (script~\texttt{20b},
multi-graines) : en VaR $99{,}5\,\%$ OpRisk le signe de l'interaction est dans le bruit de
Monte-Carlo (de $-1{,}9$ à $+2{,}1$~Md\euro{} selon la graine, $+117$~M\euro{} à haute
résolution) ; elle est en revanche \emph{super-additive} de façon robuste partout où le bruit
de queue ne domine pas : en perte moyenne sur les deux périmètres ($+528$~M\euro{} OpRisk,
$+824$~M\euro{} PRC, environ un cinquième du surcoût moyen) et jusque dans la VaR du périmètre
PRC plafonné ($+800$~M\euro, stable). Cette non-additivité motive l'allocation principielle du
surcoût (Shapley) et du niveau (Euler) du~\S\ref{sec:allocation}.

\section{Allocation du capital : Shapley et Euler}
\label{sec:allocation}

Deux questions d'allocation complètent la décomposition (script~\texttt{20}, figure S15).
\textbf{Shapley} attribue le \emph{surcoût} : fonction caractéristique
$v(S) = \mathrm{SCR}(S \text{ en NC, reste C}) - \mathrm{SCR}(\text{tous C})$ sur les $2^5$
configurations à graine commune, $\varphi_j$ la moyenne des contributions marginales de $j$
sur tous les ordres d'arrivée. Par efficience $\sum_j \varphi_j = \Delta_{\text{DORA}}$
exactement : l'interaction que les marginaux laissent de côté est intégralement redistribuée,
quel que soit son signe. Résultat (OpRisk) : P1 $34\,\%$, P4 $26\,\%$, P2 $17\,\%$,
P3 $14\,\%$, P5 $10\,\%$ du surcoût total ($10\,927$~M\euro)~; PRC $32/25/19/15/9$.
Classement $=$ ROOT sur les deux périmètres. Une part de cette convergence est mécanique
(les amorces sont semées proportionnellement à ROOT, parts $30/27/18/15/9\,\%$)~; le test non
trivial est la déformation, la part de P1 excédant sa part d'amorce ($34$ contre $30\,\%$),
effet propre de la propagation.

\textbf{Euler} alloue le \emph{niveau} $\mathrm{SCR}(\text{NC})$ par pilier touché, version
TVaR exacte : $C_j = \mathbb{E}[L_j \mid L \ge \mathrm{VaR}_{99{,}5\,\%}]$. PRC : parts
stables et presque plates (P2 $23$, P4 $23$, P1 $20$, P3 $18$, P5 $17\,\%$), capital
$\approx$ moyenne car le cap tronque la queue (TVaR $6\,795 \approx$ VaR $6\,573$~M\euro).
OpRisk : parts moyennes stables (P2 $23$, P4 $22$, P1 $20$, P3 $18$, P5 $16\,\%$), mais
l'allocation de la queue extrême n'y est pas résoluble : avec $\xi \approx 0{,}6 > 0{,}5$ la
sévérité est de variance infinie, l'estimateur de la TVaR converge en régime de loi stable et
le classement conditionnel bascule entre résolutions ($150$ et $240$~mille années). Vue
source contre vue réceptacle : Shapley concentre le surcoût là où naît la non-conformité
(P1), Euler répartit le capital là où les cascades atterrissent (profil quasi uniforme,
P2/P4 en tête légère).

\section{Markov et trajectoire du SCR(t)}
\label{sec:markov}

Le~\S\ref{sec:global} donne le SCR pour un état \emph{figé}. On ajoute la dimension temps par
une chaîne de Markov empilée sur le moteur, en deux échelles de temps (\emph{fast-slow}) : une
couche \textbf{lente} fait progresser l'état de conformité, une couche \textbf{rapide} chiffre
le SCR de chaque état par la cascade.

\subsection{La chaîne et ses séjours (2.1, 2.2)}

Chaîne à temps continu $\text{NC}\to\text{PC}\to\text{C}$, l'état Conforme étant
\textbf{absorbant} en cas de base (une fois conforme, on le reste~; la régression est un axe de
sensibilité). Les séjours ne sont pas exponentiels mais de \textbf{type phase} (Erlang-$k$) :
une exponentielle autoriserait une remédiation quasi instantanée (densité maximale en~0),
irréaliste, alors qu'un projet DORA a une durée minimale. L'Erlang-$k$ (somme de $k$
exponentielles) a un mode strictement positif, en cloche, qui capte cette durée. On développe
chaque macro-état transitoire en $k$ phases et l'on résout le générateur $Q$ sur l'espace de
phases :
\begin{equation}
\big(\mathbb P(\text{NC},t),\mathbb P(\text{PC},t),\mathbb P(\text{C},t)\big)
\;=\; \text{(sommes de phases de)}\;\; v_0\,\exp(Qt),
\end{equation}
avec $v_0$ la distribution initiale (ici la marginale ancrée $35/35/30$). Les durées moyennes
de séjour sont ancrées sur des calendriers de projet ($\sim 1{,}5$~an par transition, soit
$\sim 3$~ans de NC à~C).

\subsection{Emboîtement et couplage systémique (2.3, 2.4)}

À chaque horizon $t$, la distribution des états agrège les SCR du~\S\ref{sec:global} :
\begin{equation}
\mathrm{SCR}(t) \;=\; \mathbb P(\text{NC},t)\,\mathrm{SCR}_{\text{NC}}
+ \mathbb P(\text{PC},t)\,\mathrm{SCR}_{\text{PC}}
+ \mathbb P(\text{C},t)\,\mathrm{SCR}_{\text{C}}.
\end{equation}
Les deux couches sont \textbf{couplées par le même facteur systémique $\Theta$}, ce qui
restaure la corrélation qu'auraient perdue des couches indépendantes. Un environnement dégradé
($\Theta<0$) fait deux choses à la fois, avec le même tirage :
\begin{equation}
\underbrace{\mu(\Theta)=\mu_0\,e^{\beta_{\text{rem}}\Theta}}_{\text{ralentit la remédiation}}
\qquad\text{et}\qquad
\underbrace{\mathrm{SCR}_{\text{état}}(\Theta)=\mathrm{SCR}_{\text{état}}\,e^{-\beta_{\text{casc}}\Theta}}_{\text{durcit la cascade}}.
\end{equation}
La dispersion de $\Theta\sim\mathcal N(0,1)$ engendre la \textbf{bande d'incertitude} de la
trajectoire.

\subsection{Résultat : la trajectoire (2.5)}

Partant de l'état marginal actuel, le capital décroît vers le SCR conforme à mesure que la
remédiation aboutit, et le \textbf{Delta\_DORA($t$)}, surcoût résiduel de non-conformité, tend
vers zéro.

\begin{table}[h]\centering\small
\renewcommand{\arraystretch}{1.25}
\begin{tabular}{@{}c ccc@{}}
\toprule
$t$ (ans) & SCR($t$) médiane & bande 90\,\% & $\Delta_{\text{DORA}}(t)$ \\
\midrule
0 & 10752 & $[9092\,;12654]$ & 3827 \\
1 & 9576  & $[7205\,;12040]$ & 2652 \\
2 & 8310  & $[6194\,;11042]$ & 1385 \\
3 & 7575  & $[5924\,;10164]$ & 650 \\
5 & 7041  & $[5857\,;9028]$  & 116 \\
\bottomrule
\end{tabular}
\caption{Trajectoire du capital DORA (M\euro, OpRisk, script~\texttt{17}), vers le SCR conforme
$\approx 6925$. La bande reste large et haute : en environnement dégradé, la remédiation traîne
et le capital reste élevé plus longtemps.}
\end{table}

\section{Priorisation de la remédiation}
\label{sec:priorisation}

On ne peut pas tout remédier à la fois. Sous budget contraint (un pilier à la fois), l'ordre
de remédiation décide de la trajectoire du capital porté pendant la mise en conformité. On
cherche l'ordre qui minimise le capital intégré sur l'horizon (script~\texttt{18}).

\paragraph{Valeur d'accélération (3.1).} Gain de capital immédiat si un pilier est remédié en
premier, par différence finie depuis l'état tout-non-conforme (même logique que le tornado du
script~\texttt{15}) : $v_k = \mathrm{SCR}(\text{tous NC}) - \mathrm{SCR}(\text{tous NC sauf }k)$.
Résultat (OpRisk, base $14\,394$~M\euro) : P1 $2398$, P4 $1637$, P5 $969$, P3 $699$, P2
$\approx 0$ (dans le bruit Monte-Carlo).

\paragraph{Ordre optimal (3.2).} Un pilier remédié par période $\tau$, le coût d'un ordre est
la somme des SCR des configurations traversées,
$\text{coût}(\pi)=\tau\sum_{i=0}^{4}\mathrm{SCR}(\{\pi_1,\dots,\pi_i\})$. On énumère les
$5!=120$ ordres (SCR de chaque sous-ensemble mémoïsé).

\begin{table}[h]\centering\small
\renewcommand{\arraystretch}{1.25}
\begin{tabular}{@{}l l r@{}}
\toprule
Ordre & Séquence & capital intégré (M\euro) \\
\midrule
\textbf{Optimal} & P1 $\to$ P4 $\to$ P2 $\to$ P3 $\to$ P5 & 51471 \\
Glouton (valeur d'accél.) & P1 $\to$ P4 $\to$ P5 $\to$ P3 $\to$ P2 & 53236 \\
ROOT (qualitatif) & P1 $\to$ P4 $\to$ P2 $\to$ P3 $\to$ P5 & 51471 \\
Pire & P3 $\to$ P5 $\to$ P2 $\to$ P4 $\to$ P1 & 62273 \\
\bottomrule
\end{tabular}
\caption{Ordre de remédiation (script~\texttt{18}, OpRisk). L'écart pire/optimal est de
$21\,\%$ : l'ordre change le capital porté d'un cinquième.}
\end{table}

\begin{cle}
\textbf{Deux résultats.} D'abord, \textbf{l'ordre optimal coïncide avec l'ordre ROOT} du
modèle qualitatif : la hiérarchie posée par jugement d'expert est la séquence de remédiation
qui minimise le capital, retrouvée par une voie entièrement chiffrée. Ensuite, l'ordre optimal
\textbf{diffère de l'ordre glouton} (par valeur d'accélération immédiate) : à cause des
interactions de cascade, remédier au coup par coup le pilier au plus fort gain instantané
n'est pas optimal. Il faut raisonner sur l'horizon, pas myopiquement.
\end{cle}

\paragraph{Cohérence des approches (3.3).} La priorité de remédiation se compare à trois
autres vues : la décomposition $\Delta_k$ (\S\ref{sec:pilier}) et le classement ROOT donnent
la même hiérarchie côté amorce ($P1>P4>P2>P3>P5$) ; l'allocation d'Euler du SCR
(script~\texttt{11}, concentration de queue) remonte P4 en tête (sévérité des tiers), vue
complémentaire et non contradictoire. Trois chemins indépendants qui se recoupent sur les
piliers de tête.

\section{Robustesse et triangulation}
\label{sec:robustesse}

On teste si les verdicts survivent à la perturbation des leviers non calibrés
(script~\texttt{19}), en écho à l'analyse de sensibilité globale du modèle qualitatif.

\paragraph{Tornado du Delta\_DORA.} Base $8\,322$~M\euro. En faisant varier un levier à la
fois : l'indice de queue $\xi$ domine (de $3\,810$ à $24\,474$~M\euro, facteur $6{,}4$), comme
dans le mémoire ; la fréquence, les gains $g_{\text{NC}},g_{\text{C}}$ et le canal détection
n'ont qu'un effet modéré. \textbf{Le Delta\_DORA reste positif sur tous les réglages} : la
non-conformité coûte, quel que soit le calage.

\paragraph{Robustesse de la priorité.} Sur tous les réglages testés, \textbf{P1 reste le
pilier à remédier en premier}, et P4 le deuxième. L'ordre fin des piliers de faible impact
(P2, P3, P5, valeurs d'accélération proches) se réordonne dans le bruit Monte-Carlo, ce qui
justifie de fonder la recommandation sur l'énumération à graine commune (\S\ref{sec:priorisation})
plutôt que sur la valeur d'accélération myope.

\begin{cle}
\textbf{Verdict, et invariance.} Comme tout le chantier cascade : \textbf{le classement est
robuste, le niveau ne l'est pas} (queue lourde). Point analytique fort : l'ordre de
remédiation optimal ne dépend que des SCR par configuration, donc il est \textbf{invariant aux
taux de transition de la chaîne de Markov}, qui sont pourtant les paramètres les moins
calibrables. Ces taux fixent le \emph{calendrier} de la trajectoire, jamais la \emph{séquence}
de remédiation recommandée. La recommandation actionnable ne repose donc sur aucun paramètre
non mesurable.
\end{cle}

\begin{portee}
\textbf{Portée et limites (à énoncer devant le jury).}
\begin{itemize}[leftmargin=1.4em,itemsep=1pt,topsep=2pt]
\item \textbf{Ce qui se calibre} : la sévérité en euros (GPD sur SAS OpRisk).
\item \textbf{Ce qui est axe de sensibilité, non calibré} : le triplet de gains de propagation
  $(g_{\text{C}}<g_{\text{PC}}<g_{\text{NC}})$~; le triplet de détection~; les probabilités
  d'état et $\gamma$~; les taux de transition de la chaîne (DORA appliqué depuis 2025, aucun
  historique) et les couplages $\beta_{\text{rem}},\beta_{\text{casc}}$, ancrés sur des durées
  de projet types.
\item \textbf{En attente de validation} : le canal détection ($p_u$) se compose avec la
  fréquence et pourrait gonfler l'effet~; il est isolé en lecture séparée et exclu de la
  trajectoire tant qu'il n'est pas tranché.
\item \textbf{Données} : SAS OpRisk porte la sévérité en euros~; Data Breach n'informe que la
  structure de cascade (la direction $W$ n'est pas calibrable sur les sources ouvertes,
  cf.~script~\texttt{05}).
\end{itemize}
\end{portee}

Le script~\texttt{22} complète le tornado par les trois hypothèses restantes. Coefficients
Jacobs : $b$ stressé à $\pm 1{,}645$~SE en pivotant la droite au centroïde de la régression,
déduit des erreurs-types publiées ($\bar{x} = 10{,}04$)~; le $\xi$ dérivé va de $0{,}85$ à
$1{,}21$ mais le Delta PRC reste contenu ($3\,125$ à $3\,977$~M\euro{} autour du central
$3\,577$, le cap borne l'effet)~: axe d'hypothèse dominant du périmètre PRC, verdict intact.
$\gamma$~: le SCR espéré normal est invariant par construction ($10\,644$~M\euro, la
marginale redonne l'ancrage)~; seule la bascule en crise durcit
($\mathrm{P}(\text{NC} \mid \text{crise})$ de $84$ à $100\,\%$). Ancrage des probabilités
d'état~: le capital espéré actuel va de $9\,951$ à $11\,485$~M\euro{} selon les enquêtes
retenues, le SCR par état et le surcoût NC contre C restant inchangés (conditionnels à
l'état). Trois axes, même ligne de partage~: les hypothèses bougent des niveaux, jamais la
direction ni le classement.

\section*{Table des scripts couverts}
\addcontentsline{toc}{section}{Table des scripts}

\begin{table}[h]\centering\small
\renewcommand{\arraystretch}{1.25}
\begin{tabularx}{\linewidth}{@{}l X l@{}}
\toprule
Script & Rôle & Figure \\
\midrule
\texttt{16}  & Vasicek multi-états global : latente à seuils ordonnés, 3 canaux, SCR par état & S10 \\
\texttt{16b} & État par pilier : décomposition du Delta\_DORA, validation croisée ROOT & S11 \\
\texttt{17}  & Markov (Erlang-2) et trajectoire SCR($t$), couplage systémique $\Theta$ & S12 \\
\texttt{18}  & Priorisation de la remédiation : ordre optimal $=$ ordre ROOT & S13 \\
\texttt{19}  & Robustesse : tornado, priorité P1 robuste, invariance aux taux Markov & S14 \\
\texttt{20}  & Allocation : Shapley du surcoût ($=$ ROOT), Euler du niveau (TVaR) & S15 \\
\texttt{20b} & Interaction : signe non pinçable en VaR OpRisk, super-additive en moyenne et sur PRC & \textendash{} \\
\texttt{21}  & Résidu Jacobs : IC PRC robuste au résidu, niveau $+7\,\%$, queue dérivée épaissie & S16 \\
\texttt{22}  & Hypothèses restantes : coefficients Jacobs (pivot centroïde), $\gamma$, ancrage & S17 \\
\bottomrule
\end{tabularx}
\caption{État du chantier multi-états au moment de cette version du document.}
\end{table}

\end{document}
